\chapter*{Resumen}
\markboth{Resumen}{}
\addcontentsline{toc}{chapter}{Resumen}

Este Trabajo de Fin de Grado aborda la evolución e integración de un generador configurable de modelos de características en el entorno UVLHub, una plataforma orientada a la gestión y manipulación de modelos escritos en Universal Variability Language (UVL). La necesidad de este componente surge de la escasez de modelos UVL expresivos y configurables que permitan evaluar herramientas de análisis de variabilidad software, especialmente aquellos que incorporan niveles avanzados del lenguaje como atributos y restricciones complejas. Para ello, se ha desarrollado una arquitectura modular del generador como plugin del ecosistema Flamapy, separando la definición de los modelos, su configuración y las operaciones de generación. Además, se ha adaptado el módulo generador de UVLHub para delegar la lógica de generación en este nuevo componente independiente, mejorando la reutilización, mantenibilidad y capacidad de evolución del sistema. La solución ha sido implementada utilizando Flask (Python) y desplegada mediante contenedores Docker para garantizar un entorno reproducible. Como resultado, se ha conseguido desacoplar el generador de la aplicación principal manteniendo la funcionalidad existente y adaptando la infraestructura de pruebas asociada al módulo. El sistema cuenta con una batería de 133 pruebas automatizadas superadas y pruebas de interfaz mediante Pytest y Selenium que validan el flujo completo de generación. La principal aportación del trabajo consiste en proporcionar una base más flexible y mantenible para la generación de modelos UVL destinados a la evaluación y desarrollo de herramientas de análisis de variabilidad.

\vspace{1.5em}
\noindent\textbf{Palabras clave:} \ph{UVL, UVLHub, Flamapy, Feature Models, generación automática de modelos, Python, Docker, pruebas automatizadas.}

% --- Versión en inglés (en página aparte: \chapter* abre página nueva) ---
\chapter*{Abstract}
\markboth{Abstract}{}
\addcontentsline{toc}{chapter}{Abstract}

This project addresses the development and integration of a configurable feature model generator within the UVLHub environment, a platform designed for the management and manipulation of models written in the Universal Variability Language (UVL). The need for this component arises from the scarcity of expressive and configurable UVL models that enable the evaluation of software variability analysis tools, especially those that incorporate advanced language features such as attributes and complex constraints. To address this, a modular architecture for the generator has been developed as a plugin for the Flamapy ecosystem, separating model definition, configuration, and generation operations. In addition, the UVLHub generator module has been adapted to delegate the generation logic to this new independent component, improving the system's reusability, maintainability, and scalability. The solution was implemented using Flask (Python) and deployed via Docker containers to ensure a reproducible environment. As a result, the generator has been decoupled from the main application while maintaining existing functionality and adapting the testing infrastructure associated with the module. The system includes a suite of 133 successfully passed automated tests and interface tests using Pytest and Selenium that validate the entire generation workflow. The main contribution of this work is to provide a more flexible and maintainable foundation for generating UVL models intended for the evaluation and development of variability analysis tools.

\vspace{1.5em}
\noindent\textbf{Keywords:} \ph{UVL, UVLHub, Flamapy, Feature Models, automatic model generation, Python, Docker, automated testing.}
